Nullniveau der Lageenergie: die Wasseroberfläche. Zustand 1: am Start (nur Lageenergie). Zustand 2: im tiefsten Punkt (Lageenergie und kinetische Energie). Zustand 3: beim Eintauchen (nur kinetische Energie).

Sei $\ell$ die Seillänge und $\alpha$ der Winkel am Start. Am Start ist das Seilende $\ell\cos\alpha$ unter dem Ast, im tiefsten Punkt $\ell$ darunter. Bis zum tiefsten Punkt sinkt sie also um
\begin{align}
 \Delta h &= \hFF = \hF
\end{align}
(bei \waO\ ist das gerade die halbe Seillänge).

\begin{abcliste}
\abc
Zwischen den Zuständen 1 und 2 sinkt sie um $\Delta h$:
\begin{align}
 m\,g\,\Delta h &= \frac{1}{2}\,m\,v_1^2 \\
 v_1 &= \vaF \\
   &= \sqrt{2\cdot\g\cdot\hF} \\
   &= \va \approx \result{\vaP{0}{2}}
\end{align}

\abc
Sei $d$ die Höhe des tiefsten Punkts über dem See. Das Loslassen ändert ihre Energie nicht: Zwischen den Zuständen 1 und 3 sinkt sie um $\Delta h + d$:
\begin{align}
 m\,g\,(\Delta h+d) &= \frac{1}{2}\,m\,v_2^2 \\
 v_2 &= \vbF \\
   &= \sqrt{2\cdot\g\cdot(\hF+\dS)} \\
   &= \vb \approx \result{\vbP{0}{2}}
\end{align}
Geschwindigkeiten addieren sich nicht, Energien schon.
\end{abcliste}
